%% Complete proofs for the separately authored Part III results.

\begin{annexitem}{annex:p3.finite.1}{thm:p03-q-computably-categorical}{The rational order is computably categorical}
\begin{proof}
Let \(A\) and \(B\) be computable copies, each with domain \(\mathbb N\).
Construct an increasing sequence of finite order isomorphisms
\(f_s:A\rightharpoonup B\).  At an even stage, take the least element
\(a\notin\operatorname{dom}(f_s)\).  Its finitely many comparisons with the
current domain determine an interval of \(B\), possibly unbounded, in which
its image must lie.  Search the computable order of \(B\) for the least unused
element of that interval.  Density and the absence of endpoints guarantee
that the search terminates.  At an odd stage do the symmetric construction
for the least element outside \(\operatorname{ran}(f_s)\).  The union
\(f=\bigcup_s f_s\) is total and onto by the alternating least-element choices,
and it preserves and reflects order because each finite stage does.  To
compute \(f(a)\), simulate the construction until \(a\) enters the domain;
that stage must occur.  Hence \(f\) is a computable isomorphism.
\end{proof}
\end{annexitem}

\begin{annexitem}{annex:p3.abstract.4}{thm:p03-elementary-reduct-extension}{Extending a reduct together with its expansion}
\begin{proof}
Choose one fresh constant \(c_a\) for every element \(a\) of \(N\).
For elements of \(M\), use the same constants in the two elementary
diagrams.  Consider
\[
 D=\operatorname{ElDiag}_{\Lang}(M)
        \cup\operatorname{ElDiag}_{\Lang_0}(N)
\]
in the language \(\Lang\cup\{c_a:a\in N\}\).
The diagrams are those of Definition~\ref{def:diagrams}, with the
subscripts recording their original signatures.

Every finite subset of \(D\) is satisfiable.  To prove this, conjoin
its formulas from the second diagram into
\(\delta(\bar c_{\bar a},\bar c_{\bar b})\), where the entries of
\(\bar a\) lie in \(M\), the entries of \(\bar b\) lie in \(N\setminus M\),
and repeated constant symbols are represented by the same variable.
The tuple \(\bar b\) may be empty.
This is an \(\Lang_0\)-formula with parameters.
Since \(N\models\delta(\bar a,\bar b)\) and
\(M\mathbin{\upharpoonright}\Lang_0\preccurlyeq N\), we have
\[
 M\mathbin{\upharpoonright}\Lang_0\models\exists\bar y\,\delta(\bar a,\bar y).
\]
Choose a witnessing tuple from \(M\) and interpret the finitely many
constants for \(\bar b\) by that tuple.
Every constant naming an element of \(M\) retains its actual value.
The full \(\Lang\)-structure \(M\) then satisfies the finite fragment
of the second diagram and every selected sentence of the first.
Any remaining unused new constants can denote one fixed element of
the nonempty domain.

By compactness, \(D\) has a model \(C\).
The diagram tests, Theorem~\ref{thm:diagram-tests}, give elementary
embeddings
\[
 e_M:M\longrightarrow C\mathbin{\upharpoonright}\Lang,\qquad
 e_N:N\longrightarrow C\mathbin{\upharpoonright}\Lang_0,
\]
and \(e_N(a)=e_M(a)\) for \(a\in M\), because both maps use \(c_a^C\).
Rename the underlying set of \(C\) by a bijection sending \(e_N(a)\)
to \(a\) for every \(a\in N\), and send the other elements to a fresh
set disjoint from \(N\).  Transport all \(\Lang\)-interpretations
along this bijection and call the resulting reduct \(M^*\).
Its restriction to \(M\) agrees with the original full structure,
and its \(\Lang_0\)-restriction to \(N\) agrees with \(N\).
Both original elementary embeddings have become the required inclusions.
\end{proof}
\end{annexitem}

\begin{annexitem}{annex:p3.abstract.1}{thm:p03-robinson-joint-consistency}{Robinson joint consistency theorem}
\begin{proof}
Use first-order compactness and completeness under the ambient Choice
assumption.  In particular, consistency is equivalent to satisfiability,
and semantic consequences are derivable.

\paragraph{A complete common theory.}
For \(i=1,2\), let
\[
 C_i=\{\sigma\in\operatorname{Sent}(\Lang_0):T_i\models\sigma\},
 \qquad C=C_1\cup C_2.
\]
Every finite subset of \(C\) has a model.
Otherwise there would be finite conjunctions \(\alpha\) of members
of \(C_1\) and \(\beta\) of members of \(C_2\) such that
\(\models\beta\to\neg\alpha\).
Empty conjunctions here are \(\top\).
Then \(T_1\models\alpha\) and \(T_2\models\neg\alpha\), giving a
derivable common-language separator by completeness, contrary to
the hypothesis.
Compactness supplies a model \(A\models C\) in the common signature.
Let \(S=\operatorname{Th}_{\Lang_0}(A)\), a complete consistent theory.

For each \(i\), the union \(T_i\cup S\) is consistent.
If not, compactness would give a finite conjunction \(\delta\) of
members of \(S\) such that \(T_i\models\neg\delta\).
Thus \(\neg\delta\in C_i\subseteq S\), whereas \(\delta\in S\),
contradicting consistency of \(S\).
Choose models of both \(T_1\cup S\) and \(T_2\cup S\).
Their common reducts have complete theory \(S\).

\paragraph{Placing one reduct inside the other model.}
Start with a model \(M_0\models T_1\cup S\).
The set
\[
 T_2\cup S\cup
       \operatorname{ElDiag}_{\Lang_0}(M_0\mathbin{\upharpoonright}\Lang_0)
\]
is consistent in \(\Lang_2\) expanded by fresh constants for the
elements of \(M_0\).
Indeed, conjoin a finite diagram fragment into \(\delta(\bar c)\).
The sentence \(\exists\bar x\,\delta(\bar x)\) is true in
\(M_0\mathbin{\upharpoonright}\Lang_0\), hence belongs to \(S\).
A model of \(T_2\cup S\) therefore has witnesses for this fragment.
Interpreting its finitely many parameter constants by those witnesses
satisfies it together with the chosen theory sentences.
Compactness and the diagram test give a model \(N_0\models T_2\cup S\)
with an elementary embedding of \(M_0\mathbin{\upharpoonright}\Lang_0\) into its reduct.
Transport the structure along a bijection as in the preceding lemma
so that this embedding becomes an inclusion:
\[
 M_0\mathbin{\upharpoonright}\Lang_0\preccurlyeq N_0\mathbin{\upharpoonright}\Lang_0.
\]

\paragraph{Alternating elementary extensions.}
To bound the model choices in this recursion, fix an infinite cardinal
\(\kappa\) at least the size of \(N_0\) and of
\(\Lang_1\cup\Lang_2\).
Downward L\"owenheim--Skolem, Theorem~\ref{thm:downward-ls}, permits
each diagram model to have size at most \(\kappa\): its parameter
language has at most \(\kappa\) symbols.
Transport \(M_0,N_0\) together, placing the domain of \(M_0\) in
\(\kappa\times\{0\}\) and the remaining elements of \(N_0\) in
\(\kappa\times\{1\}\).
Place new elements at successive steps in unused blocks
\(\kappa\times\{2\},\kappa\times\{3\},\ldots\).
All subsequent choices therefore range over a fixed set of structures
on subsets of \(\kappa\times\mathbb N\); the ambient Axiom of Choice
suffices for the sequence.

Apply Lemma~\ref{thm:p03-elementary-reduct-extension} first to the
full \(\Lang_1\)-structure \(M_0\) and the common reduct of \(N_0\).
It yields a \(\Lang_1\)-structure \(M_1\) with
\[
 M_0\preccurlyeq M_1,\qquad
 N_0\mathbin{\upharpoonright}\Lang_0\preccurlyeq M_1\mathbin{\upharpoonright}\Lang_0.
\]
Apply it next to the full \(\Lang_2\)-structure \(N_0\) and the
common reduct of \(M_1\), obtaining \(N_1\) with
\[
 N_0\preccurlyeq N_1,\qquad
 M_1\mathbin{\upharpoonright}\Lang_0\preccurlyeq N_1\mathbin{\upharpoonright}\Lang_0.
\]
Continue recursively.  At each step, the transport in the lemma fixes
the entire preceding common reduct, so all these maps are literal
inclusions.  The resulting \(\Lang_1\)-structures form an elementary
chain \(M_0\preccurlyeq M_1\preccurlyeq\cdots\); the
\(\Lang_2\)-structures form an elementary chain
\(N_0\preccurlyeq N_1\preccurlyeq\cdots\).
Their common reducts interleave:
\[
 M_0\mathbin{\upharpoonright}\Lang_0\preccurlyeq N_0\mathbin{\upharpoonright}\Lang_0
 \preccurlyeq M_1\mathbin{\upharpoonright}\Lang_0\preccurlyeq N_1\mathbin{\upharpoonright}\Lang_0
 \preccurlyeq\cdots .
\]
In particular, the two chains have the same union of underlying domains.

Let \(M_\infty=\bigcup_n M_n\) and \(N_\infty=\bigcup_n N_n\).
The elementary chain theorem, Theorem~\ref{thm:elementary-chain},
gives \(M_\infty\models T_1\) and \(N_\infty\models T_2\).
Their common reducts agree.  For any finite tuple in the union domain,
choose an interleaved stage containing it; all later interpretations
of each common relation or function agree on that tuple because the
maps are embeddings.  Common constants agree for the same reason.
Consequently the union of the two compatible interpretations on this
one domain defines a \((\Lang_1\cup\Lang_2)\)-structure satisfying
\(T_1\cup T_2\), as required.

\paragraph{The complete-common-theory form.}
If both \(T_i\) extend a complete consistent \(\Lang_0\)-theory \(S\),
the no-separator hypothesis follows directly.
If \(T_1\models\sigma\) and \(T_2\models\neg\sigma\), completeness
of \(S\) puts either \(\sigma\) or \(\neg\sigma\) in \(S\).
One of the consistent \(T_i\) would then contain incompatible
consequences.  Conversely, the first part of this proof constructs
such an \(S\) compatible with both \(T_i\) from the no-separator
hypothesis.  Thus the two formulations give the same joint-consistency
principle.

\paragraph{Stone-space formulation of the first step.}
Complete \(\Lang_0\)-theories also form a Stone space with basic
clopen sets \([\sigma]=\{S:\sigma\in S\}\).
Compactness of this space is another form of first-order compactness:
if a cover by basic clopens has no finite subcover, the negations
of its defining sentences are finitely satisfiable, and a model
of all those negations gives a complete theory missed by the cover.
Every open cover refines to a cover by basic clopens.

Let \(X_i\) consist of complete common theories compatible with \(T_i\).
It is nonempty because \(T_i\) has a model, and is closed because an
incompatible complete theory has a finite conjunction witnessing
incompatibility.  The basic clopen defined by that conjunction
is disjoint from \(X_i\).
If \(X_1\) and \(X_2\) were disjoint, every point of \(X_1\) would
have a basic clopen neighborhood missing \(X_2\).
Compactness of \(X_1\) would select finitely many such neighborhoods.
Their union is one clopen \([\sigma]\), where \(\sigma\) is the
disjunction of their defining sentences.
Every \(T_1\)-model would satisfy \(\sigma\), and every
\(T_2\)-model would satisfy \(\neg\sigma\), contradicting the
no-separator hypothesis.
Hence \(X_1\cap X_2\) is nonempty, giving the same common theory
used by the elementary-chain construction.

\paragraph{Alternative model construction by ultrapowers.}
Once models \(M_i\models T_i\cup S\) have been chosen, their
\(\Lang_0\)-reducts are elementarily equivalent.
The Keisler--Shelah theorem,
Theorem~\ref{thm:p03-keisler-shelah}, supplies isomorphic ultrapowers
of those reducts.  Take the corresponding ultrapowers of the full
\(M_i\).  By \L o\'s's theorem they still satisfy \(T_i\), and
forming a reduct commutes with this ultrapower construction because
it uses the same domains, equivalence relation, and common-symbol
interpretations.
Transport the \(\Lang_2\)-interpretations along the common-reduct
isomorphism onto the domain carrying the first ultrapower.
They agree with its \(\Lang_1\)-interpretations on every common
symbol.  Their compatible union is again a joint model.
\end{proof}
\end{annexitem}

\begin{annexitem}{annex:p3.abstract.2}{thm:p03-beth-definability}{Beth definability theorem}
\begin{proof}
Let \(P'\) be a fresh predicate of the same arity \(n\) as \(P\), let
\(T'\) replace \(P\) by \(P'\) throughout \(T\), and leave every
symbol of \(\Lang\) unchanged.
Add distinct fresh constant symbols \(\bar c=(c_1,\ldots,c_n)\).
For arity zero, this tuple is empty.  If
\(T\cup\{P(\bar c)\}\) is inconsistent, then \(P\) is defined by
\(\bottom\); if \(T'\cup\{\neg P'(\bar c)\}\) is inconsistent, it is
defined by \(\neg\bottom\).
Indeed, the constants have arbitrary interpretations and do not occur
in \(T\); the first case forces \(P\) false on every tuple in every
\(T\)-model, and the second, after renaming back, forces it true.
Otherwise implicit definability makes
\[
 T\cup\{P(\bar c)\}\cup
 T'\cup\{\neg P'(\bar c)\}
\]
inconsistent: a model would give two expansions of one \(\Lang\)-reduct that
disagree on \(\bar c\).  The contrapositive of Robinson joint consistency
therefore gives a sentence \(\sigma(\bar c)\) in the common language
\(\Lang(\bar c)\) such that
\[
 T\cup\{P(\bar c)\}\entails\sigma(\bar c),
 \qquad
 T'\cup\{\neg P'(\bar c)\}\entails\neg\sigma(\bar c).
\]
The sentence \(\sigma(\bar c)\) contains neither \(P\) nor \(P'\).
The second entailment is classically equivalent to
\(T'\entails\sigma(\bar c)\to P'(\bar c)\).  Rename \(P'\) back to \(P\),
which leaves \(\sigma\) unchanged.  Together with the first entailment,
this gives \(T\models P(\bar c)\leftrightarrow\sigma(\bar c)\).

Choose distinct variables \(x_1,\ldots,x_n\), rename the bound variables
of \(\sigma\) away from them, and replace the fresh \(c_i\) by \(x_i\),
including occurrences inside terms.  Call the resulting formula
\(\theta(\bar x)\).
For every \(M\models T\) and every tuple \(\bar a\) in its domain,
interpret \(c_i\) by \(a_i\).  Substitution preserves truth under
this interpretation, so
\[
 M\models P(\bar a)\leftrightarrow\theta(\bar a).
\]
Different fresh constant symbols may denote equal elements, so this
argument covers all tuples.
Therefore
\[
 T\models\forall\bar x\,
                (P(\bar x)\leftrightarrow\theta(\bar x)),
 \qquad \FV(\theta)\subseteq\{x_1,\ldots,x_n\}.
\]
The same conclusion was established by the constant choices
\(\theta=\bot\) or \(\theta=\neg\bot\) in the two earlier cases.
It also covers an inconsistent \(T\).
For arity zero, it is simply the sentence \(P\leftrightarrow\theta\).

Apply the consequence form of first-order compactness,
Theorem~\ref{thm:fol-compactness}, to this one defining sentence.
It gives a finite \(T_0\subseteq T\) already entailing it, uniformly
for every model and tuple.
Completeness, Theorem~\ref{thm:fol-completeness}, supplies a derivation
from \(T_0\), hence from \(T\).
\end{proof}
\end{annexitem}

\begin{annexitem}{annex:p3.abstract.5}{thm:p03-craig-robinson-beth}{Craig--Robinson equivalence and Beth definability}
\begin{proof}
\paragraph{Craig interpolation implies joint consistency.}
Assume the Craig interpolation principle, and suppose
\(T_1\cup T_2\) is inconsistent.
Compactness supplies finite \(\Delta_i\subseteq T_i\) such that
\[
 \models A\to\neg B,\qquad
 A=\bigwedge\Delta_1,\quad B=\bigwedge\Delta_2.
\]
Empty conjunctions are \(\top\).
A Craig interpolant \(\sigma\) satisfies
\(\models A\to\sigma\) and \(\models\sigma\to\neg B\).
Every nonlogical symbol in \(\sigma\) occurs in both \(A\) and \(B\),
so belongs to \(\Lang_1\cap\Lang_2\).
Thus \(T_1\models\sigma\) and \(T_2\models\neg\sigma\).
Completeness converts these into the two derivable consequences
forbidden by the common-sentence hypothesis.
This proves Robinson's theorem from Craig interpolation and compactness.

\paragraph{Joint consistency implies sentence interpolation.}
Assume the common-sentence form of Robinson's theorem, and let
\(A,B\) be sentences with \(\models A\to B\).
If \(A\) is unsatisfiable, take the interpolant \(\bot\).
If \(B\) is valid, take \(\top\).
Otherwise \(\{A\}\) and \(\{\neg B\}\) are separately consistent,
and their union is inconsistent.
Give the first theory exactly the nonlogical occurrence signature of
\(A\), and the second exactly that of \(B\).
The contrapositive of Robinson's theorem supplies a sentence \(\sigma\)
in their common signature with
\[
 A\vdash\sigma,\qquad \neg B\vdash\neg\sigma.
\]
Soundness and classical contraposition yield
\(\models A\to\sigma\) and \(\models\sigma\to B\).
Logical equality and the truth constants are available regardless
of the nonlogical signature.  Hence \(\sigma\) is a Craig interpolant.

\paragraph{Free variables.}
For formulas \(A,B\), replace each variable free in either formula by
its own fresh constant symbol, obtaining sentences.
Validity under every assignment gives validity of the resulting
sentence implication under every interpretation of these constants.
The sentence argument produces an interpolant whose new constants
can come only from variables free in both original formulas.
Rename its bound variables and substitute the original variables
back for those constants.  The term and formula substitution induction
gives both original implications and the intersection condition on
free variables.  Repeated values of different parameters are allowed.

\paragraph{Beth definability.}
The proof in Annex~\ref{annex:p3.abstract.2} derives Beth definability
from Robinson's common-language separator using two renamed predicates.
Together with the first implication above, this also derives Beth
definability from Craig interpolation.
These are the implications stated in the theorem.
\end{proof}
\end{annexitem}

\begin{annexitem}{annex:p3.abstract.3}{thm:p03-lindstrom}{Lindstr\"om's theorem}
\begin{proof}
Fix an \(\mathcal L^*\)-sentence \(\alpha\) and reduce to its finite
relational occurrence signature \(\tau\).  Any constants in \(\tau\) may be
replaced definitionally by fresh unary predicates required to be singletons;
we therefore write the coding below for relation symbols only.

\emph{The finite sequence code.}
Use unary predicates \(U,V,Q_U,Q_V\), constants \(e_U,e_V\), ternary
relations \(A_U(s,a,s')\), \(A_V(t,b,t')\), and a four-place relation
\(H(s,t,a,b)\).  The intended structures have \(U,V\) as disjoint
\(\tau\)-domains, \(Q_U,Q_V\) as their finite-sequence sorts, \(e_U,e_V\) as
empty sequences, and \(A_U,A_V\) as one-coordinate append.  The graph coded
by a pair \((s,t)\) is
\[
 h_{s,t}=\{(a,b):H(s,t,a,b)\}.
\]
The following is a single first-order formula because \(\tau\) is finite:
\[
\begin{aligned}
 \operatorname{Piso}(s,t)\;\defeq{}&
  \forall a\,b\bigl(H(s,t,a,b)\to U(a)\land V(b)\bigr)\\
 &\land\forall a\,b\,b'
   \bigl(H(s,t,a,b)\land H(s,t,a,b')\to b=b'\bigr)\\
 &\land\forall a\,a'\,b
   \bigl(H(s,t,a,b)\land H(s,t,a',b)\to a=a'\bigr)\\
 &\land\bigwedge_{R\in\tau}\forall\bar a\,\bar b
   \left(\bigwedge_{j<\operatorname{arity}(R)}H(s,t,a_j,b_j)
     \to\bigl(R(\bar a)\leftrightarrow R(\bar b)\bigr)\right).
\end{aligned}
\]
Let \(\operatorname{Seq}\) be the finite conjunction saying that the four
unary sorts are pairwise disjoint, \(U,V\) are nonempty,
\(Q_U(e_U)\), \(Q_V(e_V)\), and
\[
\begin{aligned}
 &\forall a\,b\,\neg H(e_U,e_V,a,b),\\
 &\forall s\,a\bigl(Q_U(s)\land U(a)\to
      \exists s'\,(Q_U(s')\land A_U(s,a,s'))\bigr),\\
 &\forall t\,b\bigl(Q_V(t)\land V(b)\to
      \exists t'\,(Q_V(t')\land A_V(t,b,t'))\bigr),\\
 &\forall s\,a\,s'\bigl(A_U(s,a,s')\to
      Q_U(s)\land U(a)\land Q_U(s')\bigr),\\
 &\forall t\,b\,t'\bigl(A_V(t,b,t')\to
      Q_V(t)\land V(b)\land Q_V(t')\bigr),\\
 &\forall s\,t\,a\,b\,s'\,t'\bigl(
 A_U(s,a,s')\land A_V(t,b,t')\ \to{}\\[-2mm]
 &\qquad\forall x\,y\left(
 H(s',t',x,y)\leftrightarrow
   \bigl(H(s,t,x,y)\lor(x=a\land y=b)\bigr)\right)\bigr).
\end{aligned}
\]
The constants \(e_U,e_V\) are temporary constants allowed by the regular
logic's fresh-constant clause.  Only append chains starting at
\((e_U,e_V)\) will be used, so the displayed equations determine every graph
used in the argument after a standard finite number of steps.

\emph{Potential isomorphism preserves \(\mathcal L^*\)-truth.}
Add a binary relation \(I(s,t)\) and let \(\operatorname{BF}(I)\) be the
first-order conjunction
\[
\begin{aligned}
 &I(e_U,e_V),\qquad
   \forall s\,t\bigl(I(s,t)\to
      Q_U(s)\land Q_V(t)\land\operatorname{Piso}(s,t)\bigr),\\
 &\forall s\,t\,a\bigl(I(s,t)\land U(a)\to
   \exists b\,s'\,t'\,[V(b)\land A_U(s,a,s')\land A_V(t,b,t')
       \land I(s',t')]\bigr),\\
 &\forall s\,t\,b\bigl(I(s,t)\land V(b)\to
   \exists a\,s'\,t'\,[U(a)\land A_U(s,a,s')\land A_V(t,b,t')
       \land I(s',t')]\bigr).
\end{aligned}
\]
If potentially isomorphic structures \(A,B\) disagreed on \(\alpha\), first
replace a witnessing back-and-forth family by its downward closure.  It remains
a back-and-forth family, and whenever a member is extended to cover one new
element, restricting the larger map to the old graph plus that one pair gives
an exact one-pair extension still in the family.  Interpret \(I(s,t)\) on the
actual paired finite sequences exactly when their coordinate map belongs to
this downward-closed family.  The actual append operations then interpret
\(\operatorname{Seq}\land\operatorname{BF}(I)\).
Unary relativization to \(U,V\) turns the disagreement into
\(\alpha^U\land\neg\alpha^V\).  Thus the one
\(\mathcal L^*\)-sentence
\[
 \operatorname{Seq}\land\operatorname{BF}(I)
 \land\alpha^U\land\neg\alpha^V
\]
would be satisfiable.  Downward L\"owenheim--Skolem gives a countable model.
Enumerate its \(U\)- and \(V\)-sorts, start with \((e_U,e_V)\), and use the
two extension clauses alternately.  After \(n\) steps, the update axiom makes
\(h_{s_n,t_n}\) a finite map with at most \(n\) pairs, and
\(\operatorname{Piso}\) makes it a partial isomorphism.  The maps increase,
their union is total and onto, and it is therefore an isomorphism of the two
reducts.  Isomorphism invariance contradicts their disagreement on
\(\alpha\).  Hence potentially isomorphic structures agree on every
\(\mathcal L^*\)-sentence.

\emph{Bounded rank gives a first-order definition.}
Write \(A\equiv_mB\) when \(A,B\) agree on first-order sentences of
quantifier rank at most \(m\).  For fixed finite \(\tau\) and \(m\), there
are only finitely many such theories.  Prove simultaneously in \(m,n<\omega\)
that, up to logical equivalence, there are only finitely many formulas of
quantifier rank at most \(m\) whose free variables lie among the fixed list
\(x_1,\ldots,x_n\).  At \(m=0\), the finite relational signature and fixed
variable list give finitely many atoms, hence finitely many Boolean truth
functions of those atoms.  At the successor stage, every formula is equivalent
to a Boolean combination of atoms and formulas
\(\exists x_{n+1}\,\psi(x_1,\ldots,x_{n+1})\), where \(\psi\) has rank at
most \(m\); the induction hypothesis at \((m,n+1)\) again leaves only
finitely many possibilities.  Taking \(n=0\) proves the sentence claim.
Choose one finite characteristic conjunction for each
\(\equiv_m\)-class.  If \(\equiv_m\) always makes structures agree on
\(\alpha\), the disjunction of exactly the characteristic conjunctions whose
classes satisfy \(\alpha\) is a first-order sentence equivalent to
\(\alpha\).

\emph{Compactness produces potential isomorphism.}
Suppose no such \(m\) exists.  Choose \(A_m,B_m\) with
\[
 A_m\equiv_mB_m,\qquad A_m\entails\alpha,\qquad B_m\not\entails\alpha.
\]
Duplicator therefore wins the \(m\)-round EF game on \(A_m,B_m\).

Extend the finite sequence vocabulary by a unary number sort \(N\), temporary
constants \(0,c\) permitted by the fresh-constant clause, a strict order
\(<\), a binary successor relation \(S\), and a ternary relation
\(W(r,s,t)\).  Let \(\operatorname{Disc}(N,0,S,<)\) be the finite
first-order conjunction saying that \(<\) is a strict linear order on \(N\),
\(0\) is its least element, every element has a unique immediate successor,
and every nonzero element has a unique immediate predecessor, with
\[
\begin{aligned}
 &N(0)\land\forall r\,(N(r)\to\neg(r<0)),\\
 &\forall r\,(N(r)\to\exists!r'\,S(r,r')),\\
 &\forall r\,(N(r)\land r\ne0\to\exists!r'\,S(r',r)),\\
 &\forall r\,r'\left(S(r,r')\leftrightarrow
   \bigl(N(r)\land N(r')\land r<r'
     \land\neg\exists q\,(N(q)\land r<q\land q<r')\bigr)\right).
\end{aligned}
\]
Let \(\Theta\) be the conjunction of \(\operatorname{Disc}\),
\(\operatorname{Seq}\), \(N(c)\),
\(\alpha^U\land\neg\alpha^V\), and
\[
\begin{aligned}
 &W(c,e_U,e_V),\\
 &\forall r\,s\,t\bigl(W(r,s,t)\to
      N(r)\land Q_U(s)\land Q_V(t)\land\operatorname{Piso}(s,t)\bigr),\\
 &\forall r\,r'\,s\,t\bigl(S(r,r')\land W(r',s,t)\to{}\\
 &\quad [\forall a\,(U(a)\to\exists b\,s'\,t'\,
      (V(b)\land A_U(s,a,s')\land A_V(t,b,t')\land W(r,s',t')))]\\
 &\quad\land[\forall b\,(V(b)\to\exists a\,s'\,t'\,
      (U(a)\land A_U(s,a,s')\land A_V(t,b,t')\land W(r,s',t')))]\bigr).
\end{aligned}
\]
All conjuncts except the two relativized occurrences of \(\alpha\) are
first-order, so \(\Theta\) is one \(\mathcal L^*\)-sentence in the declared
finite signature of relation symbols and temporary constants.  Compactness
and downward L\"owenheim--Skolem apply uniformly to such fresh-constant
expansions by the regular logic's abstraction clause.

For each standard \(k\), use the first-order sentence
\[
 \beta_k(c)\ \defeq\
 \exists x_0\cdots x_k\left(
   x_0=0\land\bigwedge_{i<k}S(x_i,x_{i+1})\land x_k<c\right),
 \qquad
 \Gamma=\{\Theta\}\cup\{\beta_k(c):k<\omega\}.
\]
Thus no successor term has been added to the relational signature.  Every
finite subset is satisfiable: choose \(m\) larger than all its displayed chain
lengths,
interpret \(c\) by \(m\), use \(A_m,B_m\) as the \(U,V\)-reducts, and let
\(W(r,s,t)\) record the positions from which Duplicator can survive \(r\)
more rounds.  Countable compactness gives a model of \(\Gamma\), so \(c\) is
above every standard numeral.

Externally define \(r_0=c\), and let \(r_{k+1}\) be the unique predecessor of
\(r_k\).  Each \(r_k\) exists and is nonzero because the model satisfies
\(\beta_k(c)\) for every standard \(k\).  Let
\(\mathcal I_k\) contain exactly the graphs \(h_{s,t}\) for which
\(W(r_k,s,t)\) holds and there are paired append chains of length \(k\) from
\((e_U,e_V)\) to \((s,t)\).  Then \(\mathcal I_0\) contains the empty map.
The update axiom makes every member of \(\mathcal I_k\) a finite map with at
most \(k\) pairs, and the \(W\)-axiom makes it a partial isomorphism.  If
\(h_{s,t}\in\mathcal I_k\) and \(a\in U\), the forth clause produces paired
appends \(s',t'\) with \(W(r_{k+1},s',t')\); the update equation says that
\(h_{s',t'}\) extends \(h_{s,t}\) and contains \(a\) in its domain.  The back
clause gives the symmetric range extension.  Hence
\(\bigcup_{k<\omega}\mathcal I_k\) is a back-and-forth family of finite
partial isomorphisms: the two reducts are potentially isomorphic.

But \(\Theta\) says that those reducts disagree on \(\alpha\), contradicting
the potential-isomorphism result already proved.  Therefore a finite \(m\)
exists, the bounded-rank argument gives a first-order sentence equivalent to
\(\alpha\), and the arbitrary \(\mathcal L^*\)-sentence \(\alpha\) was
first-order definable.  The reverse expressive inclusion is part of the
definition of regular abstract logic.
\end{proof}
\end{annexitem}
